
% General notation
% Sem{\"e}nov arithmetic


\subsection{Basic notation}
By $\Z$ and $\N$ we denote the integers and non-negative integers,
respectively. Given an $m\times n$ integer matrix $A$, we denote by
$\norm{A}_{1,\infty}$ the $(1,\infty)$-norm of $A$, which is the
maximum over the sum of the absolute values of the coefficients of the
rows in $A$. For $\vec b\in \Z^m$, $\norm{\vec b}_\infty$ is the
largest absolute value of the numbers occurring in $\vec b$.

\subsection{Numbers as strings and strings as numbers}
Here and below, let $\Sigma=\{0,1\}$ be a binary alphabet. Any string
from $\Sigma^*$ has a natural interpretation as a binary encoding of a
natural number, possibly with an arbitrary number of leading
zeros. Conversely, any natural number in $\N$ can be converted into
its bit representation as a string in $\Sigma^*$. Finally, by
considering strings over $(\Sigma^k)^*$ for $k\ge 1$, we can represent
$k$-tuples of natural numbers as strings over $\Sigma^k$, and
\emph{vice versa}.

Formally, given $u=\vec u_n \vec u_{n-1}\dots \vec u_0 \in
(\Sigma^k)^*$, we define the tuple of natural numbers corresponding to
$u$ in \emph{most-significant digit first (msd)} notation as
%
\[
\eval{u} \defeq \sum_{i=0}^n 2^i \cdot \vec u_i\,.
\]
%
Note that $\eval{\cdot}$ is surjective but not injective.%%  Conversely,
%% for $\vec v\in \N^k \setminus \{ \vec 0 \}$, we define
%% $\ieval{n}\defeq u$ as the unique $u\in (\Sigma^k)^*$ such that
%% $\eval{u}=\vec v$ and $u\not\in \vec 0^+\cdot (\Sigma^k)^*$, i.e., the
%% msd-first encoding of $\vec v$ without any superfluous leading zeros.
%% Finally, $\ieval{\vec 0}\defeq \vec 0$.
We lift the definition of $\eval{\cdot}$
% and $\ieval{\cdot}$
to sets in the natural way.


\subsection{Generalised Sem{\"e}nov arithmetic}
For technical convenience, the structures we consider in this paper
are relational. We refer to the first-order theory of $\langle
\N,0,1,+,2^{(\cdot)} \rangle$ as \emph{Sem{\"e}nov arithmetic}, where
$+$ is the natural ternary addition relation, and $2^{(\cdot)}$ is the
power relation of base two, consisting of all tuples $(a,b) \in \N^2$
such that $b=2^a$. Sem{\"e}nov arithmetic is an extension of
Presburger arithmetic, which is the first-order theory of the
structure $\langle \N,0,1,+ \rangle$. It is known that Sem{\"e}nov
arithmetic is decidable and admits quantifier
elimination~\cite{Sem80,CP86,BCM23}.

For presentational convenience, atomic formulas of Sem{\"e}nov
arithmetic are one of the following:
%
\begin{itemize}
\item linear equations of the form $a_1 \cdot x_1 + \cdots a_d \cdot
  x_d = b$, $a_i,b \in \Z$, and
\item exponential equations of the form $x=2^y$.
\end{itemize}
%
Here, $x_1,\ldots,x_d,y$ are arbitrary first-order variables. Clearly,
richer atomic formulas such as $x + 2^{2^y} + y = z + 5$ can be
defined from those basic class of atomic formulas, since, in this
example, $x + 2^{2^y} + y = z + 5 \equiv \exists u\exists v\, u=2^v
\wedge v = 2^y \wedge x+u+y-z=5$. Moreover, since we are interpreting
numbers over non-negative integers, we can define the order relation
in existential Sem{\"e}nov arithmetic. This enables us to without loss
of generality assume that existential formulas of Sem{\"e}nov
arithmetic are positive, since $\neg (x=y) \equiv x<y \vee y<x$ and
$\neg(x=2^y) \equiv \exists z\, z=2^y \wedge \neg( x=z)$. 

The main contribution of this paper is to show that the existential
fragment of a generalisation of Sem{\"e}nov arithmetic is decidable.
Subsequently, we write $\vec 0$ to denote a tuple of 0s in any
arbitrary but fixed dimension. \emph{Generalised Sem{\"e}nov
  arithmetic} additionally allows for non-negated atomic formulas
$R(x_1,\ldots,x_k)$, where $R=\vec 0^*\cdot L$ for some regular
language $L\subseteq (\Sigma^k)^*$. We interpret $R$ as $\eval R
\subseteq \N^k$, and the additional leading zeros we require ensure
that $R=\eval{\eval{R}}^{-1}$. Subsequently, we call a language
$L\subseteq (\Sigma^k)^*$ \emph{zero closed} if $L = \vec 0^* \cdot
L$. Given a formula $\Phi(x_1,\ldots,x_n)$ of generalised Sem{\"e}nov
arithmetic, we define $\eval \Phi \subseteq \N^d$ as the set of all
satisfying assignments of $\Phi$.

The size of an atomic formula $R(x_1,\ldots,x_k)$ is defined as the
number of states of the canonical minimal DFA defining $R$. For all
other atomic formulas $\varphi$, we define their sizes $\abs \varphi$
as the number of symbols required to write down $\varphi$, assuming
binary encoding of numbers. The size $\abs \Phi$ of an arbitrary
existential formula $\Phi$ of generalised Sem{\"e}nov arithmetic is
the sum of the sizes of all atomic formulas of $\Phi$.

The full first-order theory of generalised Sem{\"e}nov arithmetic is
known to be undecidable~\cite{Point10}. This follows from the
undecidability of $\langle \N,0,1,+,2^{(\cdot)}, V_2 \rangle$, where
$V_2$ is the binary predicate such that $V_2(x,y)$ holds if and only
if $x$ is the largest power of $2$ dividing $y$ without
remainder. Note that $V_2$ can be defined in terms of a regular
language, cf.~\cite{BHMV94}. The central result of this paper is the
following:
%
\begin{theorem}\label{thm:main}
  The existential fragment of generalised Sem{\"e}nov arithmetic is
  decidable in EXPSPACE.
\end{theorem}


\subsection{Affine vector addition systems with states}
A technical tool for our decidability results is a tailor-made class
of \emph{labelled affine vector addition systems with states
  (\laVASS)}. Formally, an \laVASS is a tuple $V=\langle Q,d,
\Sigma,\Delta,\lambda,q_0,F,\Phi\rangle$, where
%
\begin{itemize}
\item $Q$ is a finite set of \emph{control states},
\item $d\ge 0$ is the \emph{dimension of $V$},
\item $\Sigma$ is a \emph{finite alphabet},
\item $\Delta \subseteq Q\times \mathcal P(\Sigma) \times Q$ is a
  finite set of \emph{transitions},
\item $\lambda\colon \Delta \to \Ops^d$ is the \emph{update function},
  where $\Ops \subseteq \Z[x]$ is the set of all affine functions over
  a single variable,
\item $q_0\in Q$ is the \emph{initial control state},
\item $F\subseteq Q$ is the set of \emph{final control states}, and
\item $\Phi$ is a a quantifier-free formula of Presburger arithmetic
  $\Phi(x_1,\ldots,x_d)$ that specifies a finite set $\eval
  \Phi\subseteq \N^d$ of \emph{final counter values}.
\end{itemize}
%
Note that when $d=0$ then $V$ is essentially a non-deterministic
finite automaton.

The set of \emph{configurations} of $V$ is $C(V)\defeq Q\times
\N^d$. The \emph{initial configuration} of $V$ is
$c_0=(q_0,0,\ldots,0)$, and the set of \emph{final configurations}
is \[C_f(V) \defeq \left \{ (q_f,\vec v) : q_f\in F,\vec v\in \eval\Phi
\right \}\,.\]
%TODO: Maybe pick some other notation than v in C_f
%
%% Given a quantifier-free formula of Presburger arithmetic $\Phi$, we
%% denote by $\abs \Phi$ the number of symbols required to write down
%% $\Phi$.
For an update function $\lambda\colon \Delta \to \Ops^d$,
we define
\[
\norm \lambda \defeq \max \{ \abs a +\abs b :
\lambda(t)=(f_1,\ldots,f_d), f_i = a x + b, 1\le i\le d, t\in T\}\,.
\]
%
We define the size $\abs V$ of an \laVASS $V=\langle Q,d,
\Sigma,\Delta,\lambda,q_0,F,S_f\rangle$ as
%
\[
\abs V \defeq \abs Q + \abs \Delta \cdot (d+1) \cdot \log (\norm \lambda + 1) +
\abs \Phi \,.
\]

An \laVASS induces an (infinite) labelled directed \emph{configuration
  graph} ${G=(C(V),\xrightarrow{}})$, where ${\xrightarrow{}}
\subseteq C(V)\times \Sigma \times C(V)$ such that $c\xrightarrow{a}
c'$ if and only if
%
\begin{itemize}
\item $c=(q,m_1,\ldots,m_d)$ and $c'=(q',m_1',\ldots,m_d')$,
\item there is $t=(q,A,q')\in \Delta$ such that
  \begin{itemize}
  \item $a\in A$,
  \item $\lambda(t)=(f_1,\ldots,f_d)$, and
  \item $m_i'=f_i(m_i)$ for all $1\le i\le d$.
  \end{itemize}
\end{itemize}
%
We lift the definition of $\xrightarrow{}$ to words $w=a_1\cdots a_n
\in \Sigma^*$ in the natural way, and thus write $c\xrightarrow{w}c'$
whenever $c \xrightarrow{a_1} c_1 \xrightarrow{a_2} \cdots c_{n-1}
\xrightarrow{a_n} c'$ for some $c_1,\ldots,c_{n-1} \in C$. The
\emph{language} $L(V) \subseteq \Sigma^*$ of $V$ is defined as
\[
L(V) \defeq \{ w \in \Sigma^* : c_0 \xrightarrow{w} c_f, c_f \in C_f(V)
\}\,.
\]
%
If we are interested in runs of \laVASS, we write $\pi = c_1\xrightarrow{t_1} c_2\xrightarrow{t_2}
\cdots \xrightarrow{t_{n-2}}c_{n-1} \xrightarrow{t_{n-1}} c_n$ to emphasise the sequence of configurations and 
transitions taken. For $1\le i\le j\le n$, we denote by $\pi[i,j]$ the
subsequence $c_i\xrightarrow{t_i} c_{i+1}\xrightarrow{t_{i+1}} \cdots \xrightarrow{t_{j-1}}c_j$. We denote by
$\val(\pi,m_i)$ the value $m_i$ of the $i$-th counter in the last configuration of $\pi$.

The \emph{emptiness problem} for an \laVASS is to decide whether
$L(V)\neq \emptyset$. Affine VASS are a powerful class of infinite
state systems, and even in the presence of only two counters and
$\Phi(x_1,x_2) \equiv x_1 = x_2$, the emptiness problem is
undecidable~\cite{Reichert2015}.  In \Cref{sec:emptiness-rlavass}, we
identify a syntactic fragment of \laVASS for which emptiness can be
decided in EXPSPACE.


\subsubsection*{Closure properties of languages of \laVASS}

We briefly discuss closure properties of \laVASS and show that they
are closed under union and intersection, and restricted kinds of
homomorphisms and inverse homomorphisms, using essentially the
standard constructions known for finite-state automata. Let
$V_i=\langle Q_i,d_i,
\Sigma,\Delta_i,\lambda_i,q_0^{(i)},F_i,\Phi_i\rangle$, $i\in
\{1,2\}$, be two \laVASS.

\begin{proposition}\label{prop:union-intersetion-closure}
  The languages of \laVASS are closed under union and intersection.
  Moreover, for $V$ such that $L(V) = L(V_1) \cap L(V_2)$, we have
  $\abs{V} \le \abs{V_1} \cdot \abs{V_2}$.
\end{proposition}
\begin{proof}
  This result can be obtained by generalising the standard
  constructions known from non-deterministic finite-state automata.
  The set of control states of the \laVASS $V$ accepting the
  intersection of \laVASS $V_1$ and $V_2$ is $Q_1\times Q_2$. The
  dimension of $V$ is the sum of the dimensions of $V_1$ and $V_2$,
  and the counters of $V_1$ and $V_2$ get independently simulated in
  the counters of $V$. Upon reading an alphabet symbol $a$, the
  \laVASS $V$ then simultaneously simulates the respective transitions
  of $V_1$ and $V_2$ for $a$; further details are relegated to
  \Cref{app:intersection}.
\end{proof}

Note that since \laVASS languages contain regular languages,
\Cref{prop:union-intersetion-closure} in particular enables us to
intersect \laVASS languages with regular languages.

Let $\Sigma,\Gamma$ be two finite alphabets. Recall that a
homomorphism $h\colon \Gamma^* \to \Sigma^*$ is fully defined by
specifying $h(a)$ for all $a\in \Gamma$. We call $h$ a
\emph{projection} if $\abs{h(a)}=1$ for all $a\in \Gamma$.
%
\begin{proposition}\label{prop:projection-closure}
  The languages of \laVASS are closed under projections and inverses
  of projections.
\end{proposition}
\begin{proof}
  Let $h\colon \Gamma^* \to \Sigma^*$ be a projection. Given an
  \laVASS $V=\langle Q,d, \Sigma,\Delta,\lambda,q_0,F,S_f\rangle$, to
  obtain closure under projections replace any $t=(q,A,q')\in \Delta$
  with $t'=(q,h(A),q')$, and set $\lambda(t')\defeq \lambda(t)$. To
  obtain closure under inverse projections, replace any $t=(q,A,q')\in
  \Delta$ with $t'=(q,h^{-1}(A),q')$ and set $\lambda(t')\defeq
  \lambda(t)$.
\end{proof}
